%%
%% Test: Layout Configuration — v1.1
%% Stage: layout refactor (stages 1-4)
%% 
%% Purpose: Verify layout settings:
%%   - Vaziri format (16.5 × 24 cm)
%%   - Asymmetric margins (top=5.3cm, bottom=5.0cm)
%%   - Header/Footer (book title / chapter name / page number)
%%   - Line spacing (1.1)
%%   - Paragraph indent (0.5cm)
%%   - draftwatermark (conditional)
%% 

\documentclass{matinbook}

\title{تست چیدمان صفحه — نسخه ۱.۱}
\author{تیم MatinBook}
\date{۱۴۰۵}

% Set book title for even-page headers
\booktitle{تست چیدمان MatinBook}

\begin{document}

\maketitle

\tableofcontents

\chapter{بررسی چیدمان صفحه}
\label{chap:test-layout}

این سند، تنظیمات چیدمان \texttt{mb-layout} را بررسی می‌کند.

%----------------------------------------
\section{قطع صفحه}
\label{sec:page-size}

\textbf{عرض صفحه:} \the\paperwidth

\textbf{ارتفاع صفحه:} \the\paperheight

\textbf{عرض متن:} \the\textwidth

\textbf{ارتفاع متن:} \the\textheight

اگر عرض صفحه حدود ۴۶۹pt (۱۶.۵cm) باشد، قطع وزیری تأیید می‌شود.

%----------------------------------------
\section{حاشیه‌ها}
\label{sec:margins}

\begin{itemize}
    \item حاشیه بالا: ۵.۳ سانتی‌متر
    \item حاشیه پایین: ۵.۰ سانتی‌متر
    \item حاشیه چپ: ۴.۵ سانتی‌متر
    \item حاشیه راست: ۴.۵ سانتی‌متر
    \item \verb|bindingoffset|: ۰.۵ سانتی‌متر
\end{itemize}

%----------------------------------------
\section{سرصفحه و پاصفحه}
\label{sec:header-footer}

در صفحات زوج، عنوان کتاب (\texttt{\textbackslash matin@booktitle}) در سرصفحه نمایش داده می‌شود.

در صفحات فرد، نام فصل (\texttt{\textbackslash leftmark}) در سرصفحه نمایش داده می‌شود.

شماره صفحه در پاصفحه قرار دارد.

%----------------------------------------
\section{فاصله خطوط}
\label{sec:line-spacing}

\textbf{فاصله خطوط:} ۱.۱

\textbf{تورفتگی پاراگراف:} ۰.۵ سانتی‌متر

این یک پاراگراف نمونه است که برای بررسی فاصله خطوط و تورفتگی نوشته شده است. در تایپوگرافی فارسی، فاصله خطوط معمولاً بین ۱.۰ تا ۱.۲ است.

%----------------------------------------
\section{محیط‌های علمی}
\label{sec:environments}

\begin{theorem}[قضیه نمونه]
    برای هر عدد حقیقی $x$، داریم $x^2 \geq 0$.
\end{theorem}

\begin{proof}
    چون مربع هر عدد حقیقی نامنفی است، نتیجه حاصل می‌شود.
\end{proof}

%----------------------------------------
\section{جدول}
\label{sec:table}

\begin{table}[h]
\centering
\caption{جدول با ستون‌های سفارشی}
\label{tab:layout-test}
\begin{tabular}{L{3cm} C{2cm} R{3cm}}
\toprule
\textbf{چپ} & \textbf{وسط} & \textbf{راست} \\
\midrule
متن & ۱۰ & ۲۰ \\
\bottomrule
\end{tabular}
\end{table}

%----------------------------------------
\section{کد}
\label{sec:code}

\begin{latin}
\begin{minted}{python}
def greet(name):
    return f"Hello, {name}!"
\end{minted}
\end{latin}

%----------------------------------------
\section{نتیجه‌گیری}
\label{sec:conclusion}

اگر این سند بدون خطا کامپایل شود:

\begin{itemize}
    \item ✅ قطع وزیری (۱۶.۵×۲۴) اعمال شده
    \item ✅ حاشیه‌های نامتقارن اعمال شده
    \item ✅ هدر/فوتر درست کار می‌کنند
    \item ✅ شماره صفحه در فوتر است
    \item ✅ فاصله خطوط ۱.۱ است
    \item ✅ تورفتگی ۰.۵ سانتی‌متر است
    \item ✅ \texttt{\textbackslash booktitle} کار می‌کند
\end{itemize}

\textbf{وضعیت تست:} قبول

\end{document}
